% Original applied explanation; see Chapter04/SOURCES.md.
\subsubsection{রুটিনের ছক থেকে বিস্তৃত ঘর গণনা}
রুটিনে একটি ল্যাব ক্লাস দুই সময়ের ঘর দখল করলে \texttt{colspan="2"} ব্যবহার করা যায়। একটি বিভাগ-নাম তিনটি বিষয়ের সারি ধরে থাকলে \texttt{rowspan="3"} উপযোগী। বিস্তৃত ঘর যত স্থান দখল করে, ওই স্থানে আর নতুন ঘর লিখবে না।

কোড লেখার আগে কাগজে সারি-কলামের ছক আঁকবে। ওপরের সারি থেকে নেমে আগের rowspan-এর দখল করা স্থান চিহ্নিত করবে; তারপর অবশিষ্ট স্থানে নতুন \texttt{td}/\texttt{th} বসাবে। এইভাবে পুরো ছকের প্রস্থ একই থাকে এবং বাড়তি ঘর তৈরির ভুল ধরা যায়।
